\section{Results and interpretation}
The objective is a fixed ordinary polynomial, not an equation with variable
exponents, a schematic proof predicate, or an uncompiled arithmetic circuit
with hidden existential wires. We count only declared unknowns. Literal
arithmetic circuits use binary addition, subtraction, and multiplication;
squaring and multiplication by a fixed numeral each cost one operation.
Sharing a previously computed register costs nothing. These are arithmetic
circuit counts, not bit-operation costs.

The central new combination removes a guard coordinate and replaces large
variable radical weights by the fixed coefficients $1,2,4,8$. A rational
projection first recovers the main root; an even shared strong ordinate
makes the other four radicands odd, so their square-class contributions
cannot cancel by $2$-adic valuation. A second new argument uses an upper
bound, rather than an absolute-value bound, on a positivity certificate to
impose exact equalities without a separate finalizer.

The later power-synchronization constructions are deliberately presented as
intermediate interfaces. Their exposed-coordinate and graph counts do not
constitute a variable count for an ordinary polynomial. The complete
polynomial results and those interfaces must not be conflated.

\subsection{Complete polynomial endpoints}
The following table separates the integer and positive domains. Degree and
operation count are independently checked after specializing the actual
fixed coefficients. These counts exclude no existential graph witnesses.
\begin{center}
\begin{tabular}{@{}lrrr@{}}
\toprule
Domain and mechanism & Variables & Degree & Operations\\
\midrule
Integers, interpolated parity, direct & 11 & 2912 & 200\\
Integers, auxiliary center & 11 & 2272 & 204\\
Integers, two centers & 11 & 2240 & 206\\
Integers, centered budget and trace & 12 & 1000 & 180\\
Integers, cheaper companion & 12 & 1024 & 179\\
Integers, bounded-guard absorption & 15 & 176 & 125\\
Strictly positive integers & 10 & 912 & 126\\
\bottomrule
\end{tabular}
\end{center}
The eleven-integer construction combines the late intersective interpolation
lemma with the centered budget and a parity separator. The twelve-integer
construction offers substantially lower degree and evaluation cost. Neither
is a variable-count record against the repository's ten-integer constructions,
which have much greater degree and use different reductions. The nine-positive
and direct quartic families likewise remain distinct comparison points.
